\documentclass[12pt,a4paper]{article}
\usepackage[utf8]{inputenc}
\usepackage[T1]{fontenc}
\usepackage[spanish]{babel}
\usepackage[margin=2.5cm]{geometry}
\usepackage{graphicx}
\usepackage{longtable}
\usepackage{array}
\usepackage{hyperref}

\graphicspath{{../03_Modelado/Diagramas_UML/}}

\title{Especificación de Requisitos de Software\\Secciones 2 y 3 --- Bloque individual}
\author{Calle Delgado Kamila Annabella}
\date{}

\begin{document}
\maketitle
\setcounter{section}{1}

% Seccion 3: Requisitos especificos completos
\section{Requisitos específicos completos}
\label{sec:requisitos-especificos}

\subsection{Requisitos funcionales}
\label{sec:requisitos-funcionales}

% Subseccion 3.1.1: Plantilla de especificacion
\subsubsection{Plantilla de especificación}
\label{sec:plantilla-rf}

Los requisitos funcionales del Sistema de Gestión Agrícola con Inteligencia Artificial para Cultivos de Verde y Cacao (SGA) se documentan mediante la plantilla de ocho atributos que se describe en la Tabla~\ref{tab:plantilla-rf}. Esta plantilla amplía la utilizada en la Entrega 2 (1B), que incluía Identificador, Nombre, Descripción, Actor(es), Prioridad, Precondiciones, Flujo principal y Postcondiciones, incorporando un origen verificable por evidencia de campo, la especificación de entradas y salidas, y un criterio de verificación cuantificable.

\begin{center}
\begin{longtable}{|p{3.5cm}|p{10.7cm}|}
\hline
\textbf{Atributo} & \textbf{Contenido} \\
\hline
\endfirsthead
\hline
\textbf{Atributo} & \textbf{Contenido} \\
\hline
\endhead
\hline
\endfoot
\hline
\endlastfoot
Identificador & Código único con el formato RF-XX. \\
\hline
Nombre & Frase nominal breve que identifica la capacidad especificada. \\
\hline
Descripción & Enunciado en la forma «El sistema deberá\ldots», limitado a una única capacidad por ficha. \\
\hline
Actor / origen & Actor que origina o utiliza el requisito, junto con el identificador de evidencia EV-XX y la marca de tiempo que lo respalda, o la restricción normativa aplicable cuando el requisito no se origina en una cita directa. \\
\hline
Entradas & Datos que el sistema recibe para ejecutar la capacidad, con su tipo y su obligatoriedad. \\
\hline
Salidas & Datos o efectos que el sistema produce como resultado de la capacidad. \\
\hline
Prioridad (MoSCoW) & Clasificación en Debe tener, Debería tener, Podría tener o No tendrá. \\
\hline
Criterio de verificación & Condición cuantificable, con métrica, umbral y método de comprobación, que permite determinar si el requisito se cumple. \\
\hline
\end{longtable}
\label{tab:plantilla-rf}
\end{center}

% Subseccion 3.1.2: Catalogo de requisitos funcionales
\subsubsection{Catálogo de requisitos funcionales}
\label{sec:catalogo-rf}

El catálogo siguiente incorpora los dieciséis requisitos funcionales identificados en la Entrega 2 (1B), migrados a la plantilla de ocho atributos sin alterar su contenido sustantivo, y su ampliación hasta un total de veintisiete requisitos funcionales derivados de la evidencia de campo descrita en la Sección 5. Los identificadores RF-01 a RF-16 corresponden a los requisitos migrados; los identificadores RF-17 a RF-27 corresponden a los requisitos incorporados en esta entrega. En los requisitos migrados se conservan, como campos adicionales al cierre de la ficha, las Precondiciones, el Flujo principal y las Postcondiciones ya redactados en la Entrega 2 (1B).

Los requisitos RF-01 y RF-14 no derivan de una cita textual de las entrevistas, ya que ninguna pregunta del guion abordó de forma directa el mecanismo de autenticación. Su origen se declara como la restricción R-02 sobre acceso diferenciado por rol, necesaria para sostener la distinción de responsabilidades entre Administrador y Trabajador agrícola evidenciada a lo largo de ambas entrevistas.

\bigskip
\noindent\textbf{RF-01. Inicio de sesión}
\label{tab:RF-01}

\begin{center}
\begin{longtable}{|p{3.5cm}|p{10.7cm}|}
\hline
Identificador & RF-01 \\
\hline
Nombre & Inicio de sesión \\
\hline
Descripción & El sistema deberá permitir que el administrador y los trabajadores agrícolas inicien sesión mediante un usuario y una contraseña para acceder a las funcionalidades correspondientes a su rol. \\
\hline
Actor / origen & Administrador, Trabajador agrícola --- R-02 (el acceso a la información dependerá del rol asignado a cada usuario). \\
\hline
Entradas & Usuario (texto, obligatorio); contraseña (texto, obligatorio). \\
\hline
Salidas & Sesión iniciada con acceso a las funciones del rol correspondiente, o mensaje de error. \\
\hline
Prioridad (MoSCoW) & Debe tener. \\
\hline
Criterio de verificación & El sistema autentica al usuario y concede acceso según su rol en $\leq$ 2,0 s, verificado en el percentil 95 de 50 intentos de inicio de sesión. \\
\hline
Precondiciones & El usuario debe estar registrado en el sistema. \\
\hline
Flujo principal & 1. El usuario ingresa sus credenciales. 2. El sistema valida la información. 3. Si las credenciales son correctas, permite el acceso al sistema según el rol del usuario. \\
\hline
Postcondiciones & El usuario accede correctamente al sistema. \\
\hline
\end{longtable}
\end{center}

\bigskip
\noindent\textbf{RF-02. Gestión de parcelas}
\label{tab:RF-02}

\begin{center}
\begin{longtable}{|p{3.5cm}|p{10.7cm}|}
\hline
Identificador & RF-02 \\
\hline
Nombre & Gestión de parcelas \\
\hline
Descripción & El sistema deberá permitir registrar, consultar, actualizar y eliminar la información correspondiente a las parcelas o lotes que conforman la finca. \\
\hline
Actor / origen & Administrador --- EV-01 [04:02], [04:14]. \\
\hline
Entradas & Nombre o código de parcela (texto, obligatorio); ubicación (texto, obligatorio); área (numérico, obligatorio); estado (selección, obligatorio). \\
\hline
Salidas & Parcela registrada, actualizada o eliminada, disponible para su asociación con cultivos (RF-03). \\
\hline
Prioridad (MoSCoW) & Debe tener. \\
\hline
Criterio de verificación & La consulta de una parcela devuelve sus campos obligatorios en $\leq$ 2,0 s con 500 parcelas cargadas, verificado en el percentil 95 de 50 ejecuciones. \\
\hline
Precondiciones & El administrador debe haber iniciado sesión. \\
\hline
Flujo principal & 1. El administrador selecciona el módulo de parcelas. 2. Registra o modifica la información requerida. 3. El sistema valida y almacena los datos. \\
\hline
Postcondiciones & La información de la parcela queda registrada o actualizada correctamente. \\
\hline
\end{longtable}
\end{center}

\bigskip
\noindent\textbf{RF-03. Gestión de cultivos}
\label{tab:RF-03}

\begin{center}
\begin{longtable}{|p{3.5cm}|p{10.7cm}|}
\hline
Identificador & RF-03 \\
\hline
Nombre & Gestión de cultivos \\
\hline
Descripción & El sistema deberá permitir registrar y administrar los cultivos de verde y cacao asociados a cada parcela de la finca. \\
\hline
Actor / origen & Administrador --- EV-01 [03:05], [03:18]. \\
\hline
Entradas & Parcela asociada (selección, obligatorio); tipo de cultivo (selección: verde o cacao, obligatorio); variedad (texto, opcional); fecha de cultivo (fecha, obligatorio). \\
\hline
Salidas & Cultivo registrado y asociado a la parcela seleccionada. \\
\hline
Prioridad (MoSCoW) & Debe tener. \\
\hline
Criterio de verificación & El registro de un cultivo queda disponible para consulta asociado a su parcela en $\leq$ 2,0 s, verificado sobre 30 registros consecutivos. \\
\hline
Precondiciones & Debe existir al menos una parcela registrada. \\
\hline
Flujo principal & 1. El administrador selecciona una parcela. 2. Registra la información del cultivo. 3. El sistema almacena los datos asociados a la parcela seleccionada. \\
\hline
Postcondiciones & El cultivo queda asociado correctamente a la parcela correspondiente. \\
\hline
\end{longtable}
\end{center}

\bigskip
\noindent\textbf{RF-04. Registro de actividades agrícolas}
\label{tab:RF-04}

\begin{center}
\begin{longtable}{|p{3.5cm}|p{10.7cm}|}
\hline
Identificador & RF-04 \\
\hline
Nombre & Registro de actividades agrícolas \\
\hline
Descripción & El sistema deberá permitir registrar las actividades realizadas por los trabajadores agrícolas, indicando la fecha, la actividad ejecutada y la parcela donde fue desarrollada. \\
\hline
Actor / origen & Trabajador agrícola --- EV-02 [00:54]; EV-01 [07:11]. \\
\hline
Entradas & Fecha (fecha, obligatorio); tipo de actividad (selección, obligatorio); parcela (selección, obligatorio). \\
\hline
Salidas & Actividad registrada, disponible para consulta del administrador (RF-10) y para el historial de trabajador (RF-24). \\
\hline
Prioridad (MoSCoW) & Debe tener. \\
\hline
Criterio de verificación & El 100~\% de las actividades registradas por un trabajador quedan disponibles para consulta del administrador en un máximo de 1 minuto, verificado sobre una muestra de 30 registros. \\
\hline
Precondiciones & El trabajador debe haber iniciado sesión. \\
\hline
Flujo principal & 1. El trabajador accede al módulo de actividades. 2. Registra la información solicitada. 3. El sistema valida y almacena el registro. \\
\hline
Postcondiciones & La actividad queda registrada correctamente en el sistema. \\
\hline
\end{longtable}
\end{center}

\bigskip
\noindent\textbf{RF-05. Registro de cosechas}
\label{tab:RF-05}

\begin{center}
\begin{longtable}{|p{3.5cm}|p{10.7cm}|}
\hline
Identificador & RF-05 \\
\hline
Nombre & Registro de cosechas \\
\hline
Descripción & El sistema deberá permitir registrar la producción obtenida en cada cosecha, especificando el cultivo, la parcela, la cantidad recolectada, la unidad de medida y la fecha del registro. \\
\hline
Actor / origen & Trabajador agrícola --- EV-01 [04:32], [04:47]; EV-02 [02:08]. \\
\hline
Entradas & Cultivo (selección, obligatorio); parcela (selección, obligatorio); cantidad (numérico, obligatorio); unidad de medida (selección, obligatorio); fecha (fecha, obligatorio). \\
\hline
Salidas & Producción registrada e incorporada al historial de producción. \\
\hline
Prioridad (MoSCoW) & Debe tener. \\
\hline
Criterio de verificación & El registro de una cosecha se incorpora al historial de producción en $\leq$ 2,0 s, verificado en el percentil 95 de 30 ejecuciones. \\
\hline
Precondiciones & Debe existir un cultivo previamente registrado en el sistema. \\
\hline
Flujo principal & 1. El trabajador selecciona el cultivo correspondiente. 2. Ingresa la información de la cosecha. 3. El sistema valida y almacena el registro. \\
\hline
Postcondiciones & La producción queda registrada y disponible para consultas y reportes. \\
\hline
\end{longtable}
\end{center}

\bigskip
\noindent\textbf{RF-06. Gestión de inventario de insumos}
\label{tab:RF-06}

\begin{center}
\begin{longtable}{|p{3.5cm}|p{10.7cm}|}
\hline
Identificador & RF-06 \\
\hline
Nombre & Gestión de inventario de insumos \\
\hline
Descripción & El sistema deberá permitir registrar, consultar, actualizar y eliminar la información correspondiente a los insumos agrícolas y herramientas utilizadas en la finca, incluyendo fertilizantes, fungicidas, herbicidas y herramientas de trabajo. \\
\hline
Actor / origen & Administrador --- EV-01 [05:36], [05:44], [05:58], [06:14]. \\
\hline
Entradas & Nombre del insumo (texto, obligatorio); tipo (selección: fertilizante, fungicida, herbicida o herramienta, obligatorio); cantidad disponible (numérico, obligatorio); unidad de medida (selección, obligatorio). \\
\hline
Salidas & Inventario actualizado, disponible para consulta y para el descuento por entregas (RF-22). \\
\hline
Prioridad (MoSCoW) & Debe tener. \\
\hline
Criterio de verificación & El inventario refleja una actualización de cantidad disponible en $\leq$ 2,0 s tras el registro, verificado sobre 30 actualizaciones consecutivas. \\
\hline
Precondiciones & El administrador debe haber iniciado sesión. \\
\hline
Flujo principal & 1. El administrador accede al módulo de inventario. 2. Registra o actualiza la información del insumo. 3. El sistema valida y almacena los cambios realizados. \\
\hline
Postcondiciones & El inventario queda actualizado con la información registrada. \\
\hline
\end{longtable}
\end{center}

\bigskip
\noindent\textbf{RF-07. Consulta de producción por período}
\label{tab:RF-07}

\begin{center}
\begin{longtable}{|p{3.5cm}|p{10.7cm}|}
\hline
Identificador & RF-07 \\
\hline
Nombre & Consulta de producción por período \\
\hline
Descripción & El sistema deberá permitir consultar la producción agrícola registrada durante un período específico, facilitando el seguimiento del rendimiento de los cultivos. \\
\hline
Actor / origen & Administrador --- EV-01 [04:47], [04:54]. \\
\hline
Entradas & Fecha de inicio y fecha de fin del período (fecha, obligatorio). \\
\hline
Salidas & Listado de producción registrada dentro del período seleccionado. \\
\hline
Prioridad (MoSCoW) & Debe tener. \\
\hline
Criterio de verificación & La consulta de producción por período con 500 registros cargados responde en $\leq$ 2,0 s, verificado en el percentil 95 de 50 ejecuciones, y presenta únicamente los registros del período seleccionado. \\
\hline
Precondiciones & Deben existir registros de producción almacenados en el sistema. \\
\hline
Flujo principal & 1. El administrador selecciona el período de consulta. 2. El sistema procesa la información registrada. 3. Se presenta el resultado de la consulta. \\
\hline
Postcondiciones & La información solicitada se muestra correctamente al usuario. \\
\hline
\end{longtable}
\end{center}

\bigskip
\noindent\textbf{RF-08. Planificación de actividades agrícolas}
\label{tab:RF-08}

\begin{center}
\begin{longtable}{|p{3.5cm}|p{10.7cm}|}
\hline
Identificador & RF-08 \\
\hline
Nombre & Planificación de actividades agrícolas \\
\hline
Descripción & El sistema deberá permitir planificar las actividades agrícolas indicando la tarea, la fecha programada y la parcela donde será ejecutada. \\
\hline
Actor / origen & Administrador --- EV-01 [06:51], [06:58]. \\
\hline
Entradas & Tarea (texto, obligatorio); fecha programada (fecha, obligatorio); parcela (selección, obligatorio). \\
\hline
Salidas & Actividad planificada, disponible para su posterior asignación (RF-09). \\
\hline
Prioridad (MoSCoW) & Debe tener. \\
\hline
Criterio de verificación & Una actividad planificada queda disponible para su asignación en $\leq$ 2,0 s desde su registro, verificado sobre 30 registros consecutivos. \\
\hline
Precondiciones & Deben existir parcelas registradas en el sistema. \\
\hline
Flujo principal & 1. El administrador accede al módulo de planificación. 2. Registra la actividad programada. 3. El sistema almacena la planificación. \\
\hline
Postcondiciones & La actividad queda programada para su posterior ejecución. \\
\hline
\end{longtable}
\end{center}

\bigskip
\noindent\textbf{RF-09. Asignación de tareas}
\label{tab:RF-09}

\begin{center}
\begin{longtable}{|p{3.5cm}|p{10.7cm}|}
\hline
Identificador & RF-09 \\
\hline
Nombre & Asignación de tareas \\
\hline
Descripción & El sistema deberá permitir asignar actividades específicas a los trabajadores agrícolas, indicando la tarea, la parcela y la fecha de ejecución. \\
\hline
Actor / origen & Administrador --- EV-01 [06:58]; EV-02 [01:01], [01:18]. \\
\hline
Entradas & Trabajador (selección, obligatorio); actividad (selección, obligatorio); parcela (selección, obligatorio); fecha de ejecución (fecha, obligatorio). \\
\hline
Salidas & Tarea asignada al trabajador seleccionado. \\
\hline
Prioridad (MoSCoW) & Debe tener. \\
\hline
Criterio de verificación & El 100~\% de las tareas asignadas quedan disponibles para consulta del trabajador correspondiente en $\leq$ 2,0 s, verificado sobre una muestra de 30 asignaciones. \\
\hline
Precondiciones & Deben existir trabajadores y actividades registrados en el sistema. \\
\hline
Flujo principal & 1. El administrador selecciona al trabajador. 2. Asigna la actividad correspondiente. 3. El sistema registra la asignación realizada. \\
\hline
Postcondiciones & La tarea queda asignada al trabajador seleccionado. \\
\hline
\end{longtable}
\end{center}

\bigskip
\noindent\textbf{RF-10. Seguimiento del cumplimiento de tareas}
\label{tab:RF-10}

\begin{center}
\begin{longtable}{|p{3.5cm}|p{10.7cm}|}
\hline
Identificador & RF-10 \\
\hline
Nombre & Seguimiento del cumplimiento de tareas \\
\hline
Descripción & El sistema deberá permitir consultar el estado de las actividades asignadas a los trabajadores, identificando aquellas pendientes, en proceso o finalizadas. \\
\hline
Actor / origen & Administrador --- EV-01 [07:11], [07:18]. \\
\hline
Entradas & Ninguna directa; opera sobre tareas previamente asignadas (RF-09). \\
\hline
Salidas & Estado de cada tarea asignada (pendiente, en proceso o finalizada). \\
\hline
Prioridad (MoSCoW) & Debe tener. \\
\hline
Criterio de verificación & El estado de una tarea asignada se actualiza y refleja correctamente en $\leq$ 2,0 s desde el cambio registrado, verificado sobre 30 tareas de prueba. \\
\hline
Precondiciones & Deben existir tareas previamente asignadas. \\
\hline
Flujo principal & 1. El administrador accede al módulo de seguimiento. 2. El sistema consulta el estado de las tareas. 3. Se presenta la información correspondiente. \\
\hline
Postcondiciones & El administrador conoce el avance de las actividades programadas. \\
\hline
\end{longtable}
\end{center}

\bigskip
\noindent\textbf{RF-11. Generación de reportes de producción}
\label{tab:RF-11}

\begin{center}
\begin{longtable}{|p{3.5cm}|p{10.7cm}|}
\hline
Identificador & RF-11 \\
\hline
Nombre & Generación de reportes de producción \\
\hline
Descripción & El sistema deberá permitir generar reportes de producción por cultivo, parcela o período, facilitando el análisis de la información registrada. \\
\hline
Actor / origen & Administrador --- EV-01 [08:10]. \\
\hline
Entradas & Filtros de consulta: cultivo, parcela o período (selección, al menos uno obligatorio). \\
\hline
Salidas & Reporte de producción según los filtros seleccionados, disponible para consulta o descarga. \\
\hline
Prioridad (MoSCoW) & Debe tener. \\
\hline
Criterio de verificación & El sistema genera el reporte de producción solicitado en $\leq$ 3,0 s con 500 registros cargados, verificado en el percentil 95 de 30 ejecuciones. \\
\hline
Precondiciones & Deben existir registros de producción almacenados en el sistema. \\
\hline
Flujo principal & 1. El administrador accede al módulo de reportes. 2. Selecciona los filtros de consulta. 3. El sistema procesa la información y genera el reporte solicitado. \\
\hline
Postcondiciones & El reporte queda disponible para su consulta o descarga. \\
\hline
\end{longtable}
\end{center}

\bigskip
\noindent\textbf{RF-12. Identificación de productividad por parcela}
\label{tab:RF-12}

\begin{center}
\begin{longtable}{|p{3.5cm}|p{10.7cm}|}
\hline
Identificador & RF-12 \\
\hline
Nombre & Identificación de productividad por parcela \\
\hline
Descripción & El sistema deberá permitir identificar las parcelas con mayor y menor productividad mediante el análisis de la producción registrada. \\
\hline
Actor / origen & Administrador --- EV-01 [07:54], [08:02]. \\
\hline
Entradas & Período de análisis (fecha, obligatorio). \\
\hline
Salidas & Listado de parcelas ordenado por producción, identificando la de mayor y la de menor productividad. \\
\hline
Prioridad (MoSCoW) & Debería tener. \\
\hline
Criterio de verificación & El sistema identifica correctamente la parcela de mayor y la de menor producción del período en el 100~\% de los casos, verificado sobre un conjunto de prueba de 20 parcelas con producción conocida. \\
\hline
Precondiciones & Deben existir registros históricos de producción por parcela. \\
\hline
Flujo principal & 1. El administrador accede al módulo de análisis. 2. El sistema procesa la información histórica. 3. Se presentan los resultados de productividad por parcela. \\
\hline
Postcondiciones & El administrador dispone de información para evaluar el rendimiento de las parcelas. \\
\hline
\end{longtable}
\end{center}

\bigskip
\noindent\textbf{RF-13. Generación de estadísticas agrícolas}
\label{tab:RF-13}

\begin{center}
\begin{longtable}{|p{3.5cm}|p{10.7cm}|}
\hline
Identificador & RF-13 \\
\hline
Nombre & Generación de estadísticas agrícolas \\
\hline
Descripción & El sistema deberá generar estadísticas basadas en la información registrada, permitiendo visualizar indicadores relacionados con la producción, las actividades realizadas y el rendimiento de los cultivos. \\
\hline
Actor / origen & Administrador --- EV-01 [07:28], [07:36], [07:43]. \\
\hline
Entradas & Ninguna directa; procesa la información ya registrada en el sistema. \\
\hline
Salidas & Gráficos e indicadores de producción, actividades y rendimiento. \\
\hline
Prioridad (MoSCoW) & Debería tener. \\
\hline
Criterio de verificación & El sistema presenta las estadísticas agrícolas en $\leq$ 3,0 s con la información de un período de 12 meses cargada, verificado en el percentil 95 de 20 ejecuciones. \\
\hline
Precondiciones & Debe existir información registrada en la base de datos. \\
\hline
Flujo principal & 1. El administrador accede al módulo de estadísticas. 2. El sistema procesa los datos disponibles. 3. Se muestran gráficos e indicadores estadísticos. \\
\hline
Postcondiciones & Las estadísticas quedan disponibles para apoyar la toma de decisiones. \\
\hline
\end{longtable}
\end{center}

\bigskip
\noindent\textbf{RF-14. Gestión de usuarios}
\label{tab:RF-14}

\begin{center}
\begin{longtable}{|p{3.5cm}|p{10.7cm}|}
\hline
Identificador & RF-14 \\
\hline
Nombre & Gestión de usuarios \\
\hline
Descripción & El sistema deberá permitir registrar, modificar, consultar y desactivar las cuentas de los usuarios autorizados, asignando el rol correspondiente a cada uno. \\
\hline
Actor / origen & Administrador --- R-02 (el acceso a la información dependerá del rol asignado a cada usuario); EV-01 [02:35]. \\
\hline
Entradas & Nombre de usuario (texto, obligatorio); rol (selección, obligatorio); estado (selección: activo o inactivo, obligatorio). \\
\hline
Salidas & Cuenta de usuario registrada, modificada o desactivada, con el rol correspondiente. \\
\hline
Prioridad (MoSCoW) & Debe tener. \\
\hline
Criterio de verificación & El 100~\% de los cambios de rol o estado de un usuario se reflejan en sus permisos de acceso en $\leq$ 2,0 s, verificado sobre una muestra de 20 cuentas de prueba. \\
\hline
Precondiciones & El administrador debe haber iniciado sesión. \\
\hline
Flujo principal & 1. El administrador accede al módulo de usuarios. 2. Registra o modifica la información del usuario. 3. El sistema valida los datos y guarda los cambios. \\
\hline
Postcondiciones & La información del usuario queda actualizada en el sistema. \\
\hline
\end{longtable}
\end{center}

\bigskip
\noindent\textbf{RF-15. Consultar recomendaciones de IA}
\label{tab:RF-15}

\begin{center}
\begin{longtable}{|p{3.5cm}|p{10.7cm}|}
\hline
Identificador & RF-15 \\
\hline
Nombre & Consultar recomendaciones de IA \\
\hline
Descripción & El sistema deberá analizar la información registrada sobre los cultivos de verde y cacao para generar recomendaciones que apoyen la planificación de actividades y la toma de decisiones por parte del administrador. \\
\hline
Actor / origen & Administrador --- EV-01 [09:26], [09:36]. \\
\hline
Entradas & Ninguna directa; el módulo de IA procesa los registros históricos existentes de producción, actividades e inventario. \\
\hline
Salidas & Recomendaciones relacionadas con la gestión de los cultivos, sujetas a la restricción R-05. \\
\hline
Prioridad (MoSCoW) & Podría tener. \\
\hline
Criterio de verificación & El sistema presenta recomendaciones fundamentadas en la información registrada en $\leq$ 5,0 s, sin modificar automáticamente los datos almacenados, verificado sobre 20 solicitudes de recomendación. \\
\hline
Precondiciones & Deben existir registros históricos suficientes de producción, actividades e inventario. \\
\hline
Flujo principal & 1. El administrador accede al módulo de Inteligencia Artificial. 2. El sistema analiza la información disponible. 3. Se generan recomendaciones relacionadas con la gestión de los cultivos. \\
\hline
Postcondiciones & Las recomendaciones quedan disponibles para su consulta por parte del administrador. \\
\hline
\end{longtable}
\end{center}

\bigskip
\noindent\textbf{RF-16. Registro de enfermedades del cultivo}
\label{tab:RF-16}

\begin{center}
\begin{longtable}{|p{3.5cm}|p{10.7cm}|}
\hline
Identificador & RF-16 \\
\hline
Nombre & Registro de enfermedades del cultivo \\
\hline
Descripción & El sistema deberá permitir registrar las enfermedades detectadas en los cultivos de verde y cacao, indicando el cultivo afectado, el nombre de la enfermedad, el nivel de gravedad y observaciones, con el fin de llevar un control fitosanitario y facilitar la generación de alertas cuando sea necesario. \\
\hline
Actor / origen & Administrador, Trabajador agrícola --- EV-02 [01:42]; EV-01 [01:16], [01:25]. \\
\hline
Entradas & Cultivo afectado (selección, obligatorio); nombre de la enfermedad (selección desde catálogo, RF-20, obligatorio); nivel de gravedad (selección, obligatorio); observaciones (texto, opcional). \\
\hline
Salidas & Enfermedad registrada y, si la gravedad es alta, alerta generada (RF-25). \\
\hline
Prioridad (MoSCoW) & Debería tener. \\
\hline
Criterio de verificación & El 100~\% de los registros de enfermedad con nivel de gravedad alto generan una alerta en $\leq$ 5,0 s desde su registro, verificado sobre una muestra de 30 casos. \\
\hline
Precondiciones & Debe existir un cultivo previamente registrado en el sistema. \\
\hline
Flujo principal & 1. El usuario selecciona el cultivo afectado. 2. Registra la enfermedad detectada. 3. Ingresa el nivel de gravedad y observaciones. 4. El sistema valida y almacena la información. 5. Si la gravedad es alta, el sistema genera una alerta. \\
\hline
Postcondiciones & La enfermedad queda registrada y disponible para su consulta. \\
\hline
\end{longtable}
\end{center}

A continuación se incorporan los once requisitos funcionales adicionales (RF-17 a RF-27), derivados de las capacidades identificadas en la evidencia de campo que no quedaban cubiertas por los dieciséis requisitos migrados. Cada uno sigue estrictamente la plantilla de ocho atributos, sin campos adicionales.

\bigskip
\noindent\textbf{RF-17. Registro de peso y comprobante de acopio}
\label{tab:RF-17}

\begin{center}
\begin{longtable}{|p{3.5cm}|p{10.7cm}|}
\hline
Identificador & RF-17 \\
\hline
Nombre & Registro de peso y comprobante de acopio \\
\hline
Descripción & El sistema deberá permitir registrar el peso entregado y el número de comprobante emitido por el centro de acopio para cada cosecha, como respaldo de la producción declarada. \\
\hline
Actor / origen & Trabajador agrícola, Administrador --- EV-01 [04:32], [04:47]. \\
\hline
Entradas & Peso entregado (numérico, kg o toneladas, obligatorio); número de comprobante (texto, obligatorio). \\
\hline
Salidas & Registro de cosecha (RF-05) con peso y comprobante asociados. \\
\hline
Prioridad (MoSCoW) & Debería tener. \\
\hline
Criterio de verificación & El 100~\% de los registros de cosecha con comprobante de acopio permiten recuperar el peso y el número de comprobante en una consulta posterior, verificado sobre una muestra de 30 registros. \\
\hline
\end{longtable}
\end{center}

\bigskip
\noindent\textbf{RF-18. Comparación de producción entre periodos}
\label{tab:RF-18}

\begin{center}
\begin{longtable}{|p{3.5cm}|p{10.7cm}|}
\hline
Identificador & RF-18 \\
\hline
Nombre & Comparación de producción entre periodos \\
\hline
Descripción & El sistema deberá permitir comparar la producción registrada entre dos periodos seleccionados por el administrador, para un mismo cultivo o parcela. \\
\hline
Actor / origen & Administrador --- EV-01 [04:47], [04:54]. \\
\hline
Entradas & Periodo 1, periodo 2, cultivo o parcela (selección, obligatorio). \\
\hline
Salidas & Diferencia porcentual y absoluta de producción entre los dos periodos seleccionados. \\
\hline
Prioridad (MoSCoW) & Debería tener. \\
\hline
Criterio de verificación & La comparación entre dos periodos con 500 registros de producción se presenta en $\leq$ 3,0 s, verificado en el percentil 95 de 30 ejecuciones. \\
\hline
\end{longtable}
\end{center}

\bigskip
\noindent\textbf{RF-19. Reporte de pérdidas por plagas}
\label{tab:RF-19}

\begin{center}
\begin{longtable}{|p{3.5cm}|p{10.7cm}|}
\hline
Identificador & RF-19 \\
\hline
Nombre & Reporte de pérdidas por plagas \\
\hline
Descripción & El sistema deberá generar un reporte de pérdidas de producción estimadas asociadas a plagas y enfermedades, por cultivo y por parcela. \\
\hline
Actor / origen & Administrador --- EV-01 [08:17], [08:22], [08:29] (reporte señalado como el más importante por el administrador). \\
\hline
Entradas & Periodo (fecha, obligatorio); cultivo o parcela (selección, opcional). \\
\hline
Salidas & Reporte con pérdida estimada por plaga o enfermedad, cultivo y parcela. \\
\hline
Prioridad (MoSCoW) & Debe tener. \\
\hline
Criterio de verificación & El reporte de pérdidas por plagas se genera en $\leq$ 3,0 s con 500 registros de enfermedad cargados, verificado en el percentil 95 de 30 ejecuciones. \\
\hline
\end{longtable}
\end{center}

\bigskip
\noindent\textbf{RF-20. Catálogo de plagas y enfermedades por cultivo}
\label{tab:RF-20}

\begin{center}
\begin{longtable}{|p{3.5cm}|p{10.7cm}|}
\hline
Identificador & RF-20 \\
\hline
Nombre & Catálogo de plagas y enfermedades por cultivo \\
\hline
Descripción & El sistema deberá mantener un catálogo de plagas y enfermedades asociadas a cada tipo de cultivo, con nombre normalizado, para estandarizar el registro de incidencias fitosanitarias. \\
\hline
Actor / origen & Administrador --- EV-02 [01:42]; EV-01 [01:16], [01:25]. \\
\hline
Entradas & Nombre de la plaga o enfermedad (texto, obligatorio); cultivo asociado (selección, obligatorio). \\
\hline
Salidas & Lista de plagas y enfermedades disponible para selección en el registro de enfermedades (RF-16). \\
\hline
Prioridad (MoSCoW) & Debería tener. \\
\hline
Criterio de verificación & El 100~\% de los registros de enfermedad (RF-16) permiten seleccionar el nombre desde el catálogo sin ingreso de texto libre, verificado sobre una muestra de 30 registros. \\
\hline
\end{longtable}
\end{center}

\bigskip
\noindent\textbf{RF-21. Planificación de compra de insumos por calendario}
\label{tab:RF-21}

\begin{center}
\begin{longtable}{|p{3.5cm}|p{10.7cm}|}
\hline
Identificador & RF-21 \\
\hline
Nombre & Planificación de compra de insumos por calendario \\
\hline
Descripción & El sistema deberá permitir planificar la compra de insumos agrícolas asignando cantidades estimadas por cultivo a meses específicos del año. \\
\hline
Actor / origen & Administrador --- EV-01 [05:36], [05:44]. \\
\hline
Entradas & Insumo (selección, obligatorio); cultivo (selección, obligatorio); mes (selección, obligatorio); cantidad estimada (numérico, obligatorio). \\
\hline
Salidas & Calendario de compras planificadas por insumo y mes. \\
\hline
Prioridad (MoSCoW) & Debería tener. \\
\hline
Criterio de verificación & El sistema muestra el calendario de compras planificadas de los 12 meses en $\leq$ 2,0 s, verificado en el percentil 95 de 20 ejecuciones. \\
\hline
\end{longtable}
\end{center}

\bigskip
\noindent\textbf{RF-22. Registro de entrega de insumos a trabajadores}
\label{tab:RF-22}

\begin{center}
\begin{longtable}{|p{3.5cm}|p{10.7cm}|}
\hline
Identificador & RF-22 \\
\hline
Nombre & Registro de entrega de insumos a trabajadores \\
\hline
Descripción & El sistema deberá permitir registrar la entrega de insumos agrícolas del administrador a un trabajador, actualizando la disponibilidad del inventario. \\
\hline
Actor / origen & Administrador --- EV-02 [02:25]. \\
\hline
Entradas & Insumo (selección, obligatorio); cantidad entregada (numérico, obligatorio); trabajador receptor (selección, obligatorio); fecha (fecha, obligatorio). \\
\hline
Salidas & Actualización del inventario (RF-06) y registro de entrega asociado al trabajador. \\
\hline
Prioridad (MoSCoW) & Podría tener. \\
\hline
Criterio de verificación & El 100~\% de las entregas registradas descuentan correctamente la cantidad del inventario disponible, verificado sobre una muestra de 30 entregas. \\
\hline
\end{longtable}
\end{center}

\bigskip
\noindent\textbf{RF-23. Registro de ingresos}
\label{tab:RF-23}

\begin{center}
\begin{longtable}{|p{3.5cm}|p{10.7cm}|}
\hline
Identificador & RF-23 \\
\hline
Nombre & Registro de ingresos \\
\hline
Descripción & El sistema deberá permitir registrar los ingresos generados por la venta de la producción agrícola, asociados a cultivo y periodo. \\
\hline
Actor / origen & Administrador --- EV-01 [06:23], [06:28], [06:37]. \\
\hline
Entradas & Monto (numérico, obligatorio); cultivo (selección, obligatorio); periodo (fecha, obligatorio); fecha (fecha, obligatorio). \\
\hline
Salidas & Registro de ingreso disponible para consulta por periodo. \\
\hline
Prioridad (MoSCoW) & Podría tener (necesidad secundaria frente a las prioridades explícitas del administrador). \\
\hline
Criterio de verificación & El registro de un ingreso queda disponible para consulta en $\leq$ 2,0 s tras su ingreso, verificado sobre 20 registros consecutivos. \\
\hline
\end{longtable}
\end{center}

\bigskip
\noindent\textbf{RF-24. Consulta de historial de registros por trabajador y fecha}
\label{tab:RF-24}

\begin{center}
\begin{longtable}{|p{3.5cm}|p{10.7cm}|}
\hline
Identificador & RF-24 \\
\hline
Nombre & Consulta de historial de registros por trabajador y fecha \\
\hline
Descripción & El sistema deberá permitir consultar el historial de labores y cosechas registradas, filtrado por trabajador y por rango de fechas, para permitir la reconciliación de la información y evitar la desviación de datos identificada en el registro en papel. \\
\hline
Actor / origen & Administrador --- EV-01 [02:04], [02:17]. \\
\hline
Entradas & Trabajador (selección, opcional); rango de fechas (obligatorio). \\
\hline
Salidas & Listado de labores y cosechas registradas en el rango filtrado. \\
\hline
Prioridad (MoSCoW) & Debe tener. \\
\hline
Criterio de verificación & La consulta del historial filtrado por trabajador y fecha con 500 registros se presenta en $\leq$ 2,0 s, verificado en el percentil 95 de 30 ejecuciones. \\
\hline
\end{longtable}
\end{center}

\bigskip
\noindent\textbf{RF-25. Generación de alertas automáticas de atención inmediata}
\label{tab:RF-25}

\begin{center}
\begin{longtable}{|p{3.5cm}|p{10.7cm}|}
\hline
Identificador & RF-25 \\
\hline
Nombre & Generación de alertas automáticas de atención inmediata \\
\hline
Descripción & El sistema deberá generar alertas automáticas visibles al administrador cuando se registre una condición que requiera atención inmediata, incluyendo enfermedades de gravedad alta e insumos con disponibilidad crítica. \\
\hline
Actor / origen & Administrador --- EV-01 [10:33]. \\
\hline
Entradas & Ninguna directa; se genera a partir de los registros de RF-16 y RF-06. \\
\hline
Salidas & Alerta con tipo, origen (parcela o insumo) y momento de generación. \\
\hline
Prioridad (MoSCoW) & Debe tener. \\
\hline
Criterio de verificación & El 100~\% de los registros de enfermedad con gravedad alta generan una alerta visible en $\leq$ 5,0 s desde su registro, verificado sobre una muestra de 30 casos. \\
\hline
\end{longtable}
\end{center}

\bigskip
\noindent\textbf{RF-26. Detección temprana de señales de enfermedad en parcelas mediante IA}
\label{tab:RF-26}

\begin{center}
\begin{longtable}{|p{3.5cm}|p{10.7cm}|}
\hline
Identificador & RF-26 \\
\hline
Nombre & Detección temprana de señales de enfermedad en parcelas mediante IA \\
\hline
Descripción & El sistema deberá analizar mediante IA la información histórica registrada de una parcela para generar una alerta temprana cuando se identifiquen patrones asociados a una posible enfermedad, antes de su confirmación visual por el usuario. \\
\hline
Actor / origen & Administrador --- EV-01 [09:02]. \\
\hline
Entradas & Registros históricos de la parcela (labores, cosechas, enfermedades previas). \\
\hline
Salidas & Alerta temprana de posible enfermedad, con parcela y nivel de confianza del modelo. \\
\hline
Prioridad (MoSCoW) & Debe tener. \\
\hline
Criterio de verificación & El módulo de IA genera la alerta temprana con una exactitud (accuracy) mínima del 70~\% sobre un conjunto de validación de al menos 30 casos etiquetados. \\
\hline
\end{longtable}
\end{center}

\bigskip
\noindent\textbf{RF-27. Consulta de producción agregada por cultivo}
\label{tab:RF-27}

\begin{center}
\begin{longtable}{|p{3.5cm}|p{10.7cm}|}
\hline
Identificador & RF-27 \\
\hline
Nombre & Consulta de producción agregada por cultivo \\
\hline
Descripción & El sistema deberá permitir consultar la producción total registrada agrupada por cultivo, sin necesidad de discriminar por parcela individual. \\
\hline
Actor / origen & Administrador --- EV-01 [09:57]. \\
\hline
Entradas & Cultivo (selección); periodo (opcional). \\
\hline
Salidas & Producción total agregada del cultivo seleccionado. \\
\hline
Prioridad (MoSCoW) & Debería tener. \\
\hline
Criterio de verificación & La consulta de producción agregada por cultivo con 500 registros se presenta en $\leq$ 2,0 s, verificado en el percentil 95 de 30 ejecuciones. \\
\hline
\end{longtable}
\end{center}

\subsection{Requisitos no funcionales}
\label{sec:requisitos-no-funcionales}

% Subseccion 3.2.1: Criterio de cuantificacion ISO/IEC 25010:2023
\subsubsection{Criterio de cuantificación ISO/IEC 25010:2023}
\label{sec:criterio-cuantificacion-rnf}

La cuantificación de los requisitos no funcionales del SGA sigue un criterio único: todo requisito no funcional debe expresarse mediante una métrica, un valor objetivo y un método de medición verificable, evitando criterios procedimentales que solo describan una actividad de prueba sin un umbral de aceptación. Este criterio se aplica sobre las nueve características del modelo de calidad de producto de la norma ISO/IEC 25010:2023~\cite{iso25010-2023}. La Tabla~\ref{tab:cobertura-rnf} resume la cobertura mínima exigida por característica y la cobertura alcanzada en el catálogo de la Sección~\ref{sec:catalogo-rnf}.

\begin{center}
\begin{longtable}{|p{5cm}|p{2.2cm}|p{6cm}|}
\hline
\textbf{Característica} & \textbf{Mínimo} & \textbf{RNF incluidos} \\
\hline
\endfirsthead
\hline
\textbf{Característica} & \textbf{Mínimo} & \textbf{RNF incluidos} \\
\hline
\endhead
\hline
\endfoot
\hline
\endlastfoot
Adecuación funcional & 1 & RNF-07 \\
\hline
Eficiencia de desempeño & 2 & RNF-01, RNF-08 \\
\hline
Compatibilidad & 1 & RNF-05 \\
\hline
Interacción de usuario & 2 & RNF-04, RNF-12 \\
\hline
Fiabilidad & 2 & RNF-02, RNF-09 \\
\hline
Seguridad & 3 & RNF-03, RNF-10, RNF-11 \\
\hline
Mantenibilidad & 1 & RNF-06 \\
\hline
Portabilidad & 1 & RNF-13 \\
\hline
Flexibilidad & 1 & RNF-14 \\
\hline
Explicabilidad (Sección~\ref{sec:explicabilidad-ia}) & 1 & RNF-15 \\
\hline
\end{longtable}
\label{tab:cobertura-rnf}
\end{center}

% Subseccion 3.2.2: Catalogo de requisitos no funcionales
\subsubsection{Catálogo de requisitos no funcionales}
\label{sec:catalogo-rnf}

Cada requisito no funcional se documenta mediante la ficha de diez atributos descrita en la Tabla~\ref{tab:plantilla-rnf}.

\begin{center}
\begin{longtable}{|p{3.8cm}|p{10.4cm}|}
\hline
\textbf{Atributo} & \textbf{Contenido} \\
\hline
\endfirsthead
\hline
\textbf{Atributo} & \textbf{Contenido} \\
\hline
\endhead
\hline
\endfoot
\hline
\endlastfoot
Identificador & Código único con el formato RNF-XX. \\
\hline
Nombre & Frase nominal breve que identifica la propiedad de calidad especificada. \\
\hline
Característica ISO/IEC 25010:2023 & Una de las nueve características del modelo de calidad de producto. \\
\hline
Subcaracterística & Subcaracterística específica dentro de la característica declarada. \\
\hline
Descripción & Enunciado en la forma «El sistema deberá\ldots», referido a una propiedad de calidad medible. \\
\hline
Métrica & Variable que se mide para verificar el cumplimiento del requisito. \\
\hline
Valor objetivo & Umbral numérico o condición que la métrica debe alcanzar. \\
\hline
Método de medición & Procedimiento, tamaño de muestra y condiciones bajo las cuales se mide la métrica. \\
\hline
Origen & Identificador de evidencia EV-XX con marca de tiempo, o referencia normativa cuando el requisito no se origina en una cita directa. \\
\hline
Prioridad & Alta, Media o Baja. \\
\hline
\end{longtable}
\label{tab:plantilla-rnf}
\end{center}

Los requisitos RNF-01 a RNF-06 corresponden a los identificados en la Entrega 2 (1B), migrados a esta ficha y corregidos: sus criterios de verificación originales describían una actividad de prueba sin un umbral cuantificable, y se sustituyen aquí por una métrica con valor objetivo y método de medición explícitos. La categoría «Usabilidad» de la Entrega 2 (1B) se renombra a «Interacción de usuario», conforme a la nomenclatura de la edición 2023 de la norma. Los requisitos RNF-07 a RNF-14 se incorporan para cubrir las características y los mínimos exigidos en esta entrega. El requisito RNF-15 se desarrolla íntegramente en la Sección~\ref{sec:explicabilidad-ia}.

\bigskip
\noindent\textbf{RNF-01. Tiempo de respuesta en la consulta de parcelas}
\label{tab:RNF-01}

\begin{center}
\begin{longtable}{|p{4.2cm}|p{10cm}|}
\hline
Identificador & RNF-01 \\
\hline
Nombre & Tiempo de respuesta en la consulta de parcelas \\
\hline
Característica ISO/IEC 25010:2023 & Eficiencia de desempeño \\
\hline
Subcaracterística & Comportamiento temporal \\
\hline
Descripción & El sistema deberá responder a las consultas y operaciones del usuario dentro de un tiempo máximo bajo condiciones normales de uso. \\
\hline
Métrica & Tiempo de respuesta \\
\hline
Valor objetivo & El listado de parcelas responde en $\leq$ 2,0 s con 500 registros cargados. \\
\hline
Método de medición & Medición del percentil 95 de 50 ejecuciones. \\
\hline
Origen & ISO/IEC 25010:2023 (característica sin evidencia directa de campo); RF-02. \\
\hline
Prioridad & Alta \\
\hline
\end{longtable}
\end{center}

\bigskip
\noindent\textbf{RNF-02. Disponibilidad del sistema}
\label{tab:RNF-02}

\begin{center}
\begin{longtable}{|p{4.2cm}|p{10cm}|}
\hline
Identificador & RNF-02 \\
\hline
Nombre & Disponibilidad del sistema \\
\hline
Característica ISO/IEC 25010:2023 & Fiabilidad \\
\hline
Subcaracterística & Disponibilidad \\
\hline
Descripción & El sistema deberá estar disponible durante el horario de operación de la finca. \\
\hline
Métrica & Porcentaje de tiempo disponible (uptime) \\
\hline
Valor objetivo & Disponibilidad $\geq$ 99,0~\% mensual en horario 06:00--18:00. \\
\hline
Método de medición & Medida como uptime registrado durante 30 días. \\
\hline
Origen & ISO/IEC 25010:2023 (característica sin evidencia directa de campo). \\
\hline
Prioridad & Alta \\
\hline
\end{longtable}
\end{center}

\bigskip
\noindent\textbf{RNF-03. Control de acceso por rol}
\label{tab:RNF-03}

\begin{center}
\begin{longtable}{|p{4.2cm}|p{10cm}|}
\hline
Identificador & RNF-03 \\
\hline
Nombre & Control de acceso por rol \\
\hline
Característica ISO/IEC 25010:2023 & Seguridad \\
\hline
Subcaracterística & Autenticidad \\
\hline
Descripción & El sistema deberá permitir el acceso únicamente a usuarios autenticados mediante credenciales válidas y controlar los permisos según el rol asignado. \\
\hline
Métrica & Porcentaje de intentos de acceso fuera de rol rechazados \\
\hline
Valor objetivo & El 100~\% de los intentos de acceso a funciones fuera del rol asignado son rechazados. \\
\hline
Método de medición & Verificado sobre una matriz de 12 casos rol-función. \\
\hline
Origen & R-02, RD-04. \\
\hline
Prioridad & Alta \\
\hline
\end{longtable}
\end{center}

\bigskip
\noindent\textbf{RNF-04. Curva de aprendizaje de usuarios nuevos}
\label{tab:RNF-04}

\begin{center}
\begin{longtable}{|p{4.2cm}|p{10cm}|}
\hline
Identificador & RNF-04 \\
\hline
Nombre & Curva de aprendizaje de usuarios nuevos \\
\hline
Característica ISO/IEC 25010:2023 & Interacción de usuario \\
\hline
Subcaracterística & Capacidad de aprendizaje \\
\hline
Descripción & Las interfaces del sistema deberán ser suficientemente intuitivas para que un usuario nuevo realice las funciones básicas tras una capacitación breve. \\
\hline
Métrica & Porcentaje de usuarios de prueba que completan las funciones básicas tras la capacitación \\
\hline
Valor objetivo & El 100~\% de los usuarios de prueba completan las funciones básicas tras una capacitación de 30 minutos. \\
\hline
Método de medición & Prueba con al menos 5 usuarios nuevos, midiendo el tiempo necesario para ejecutar las tareas principales. \\
\hline
Origen & ISO/IEC 25010:2023, motivado por el perfil de usuario sin experiencia digital previa declarado en EV-02 [00:24]. \\
\hline
Prioridad & Media \\
\hline
\end{longtable}
\end{center}

\bigskip
\noindent\textbf{RNF-05. Compatibilidad con navegadores}
\label{tab:RNF-05}

\begin{center}
\begin{longtable}{|p{4.2cm}|p{10cm}|}
\hline
Identificador & RNF-05 \\
\hline
Nombre & Compatibilidad con navegadores \\
\hline
Característica ISO/IEC 25010:2023 & Compatibilidad \\
\hline
Subcaracterística & Coexistencia \\
\hline
Descripción & El sistema deberá funcionar correctamente en los navegadores Google Chrome, Microsoft Edge y Mozilla Firefox en sus versiones actuales. \\
\hline
Métrica & Proporción de casos de uso ejecutados sin error funcional, por navegador \\
\hline
Valor objetivo & 18 de 18 ejecuciones exitosas (6 casos de uso $\times$ 3 navegadores), 0 errores funcionales. \\
\hline
Método de medición & Ejecución de los seis casos de uso principales (CU-01 a CU-06 de la Entrega 2, 1B) en cada uno de los tres navegadores compatibles. \\
\hline
Origen & RD-03. \\
\hline
Prioridad & Media \\
\hline
\end{longtable}
\end{center}

\bigskip
\noindent\textbf{RNF-06. Modularidad de la arquitectura}
\label{tab:RNF-06}

\begin{center}
\begin{longtable}{|p{4.2cm}|p{10cm}|}
\hline
Identificador & RNF-06 \\
\hline
Nombre & Modularidad de la arquitectura \\
\hline
Característica ISO/IEC 25010:2023 & Mantenibilidad \\
\hline
Subcaracterística & Modularidad \\
\hline
Descripción & El sistema deberá estar desarrollado utilizando una arquitectura modular que facilite la incorporación de nuevas funcionalidades y la corrección de errores. \\
\hline
Métrica & Número de módulos modificados al incorporar una funcionalidad de complejidad equivalente a un módulo existente \\
\hline
Valor objetivo & Una nueva funcionalidad (por ejemplo, un nuevo tipo de reporte) se incorpora modificando como máximo 3 módulos del sistema. \\
\hline
Método de medición & Revisión de la arquitectura documentada por un desarrollador independiente al cambio. \\
\hline
Origen & RD-05. \\
\hline
Prioridad & Media \\
\hline
\end{longtable}
\end{center}

\bigskip
\noindent\textbf{RNF-07. Cobertura de unidades de medida por cultivo}
\label{tab:RNF-07}

\begin{center}
\begin{longtable}{|p{4.2cm}|p{10cm}|}
\hline
Identificador & RNF-07 \\
\hline
Nombre & Cobertura de unidades de medida por cultivo \\
\hline
Característica ISO/IEC 25010:2023 & Adecuación funcional \\
\hline
Subcaracterística & Completitud funcional \\
\hline
Descripción & El sistema deberá permitir registrar la cosecha en la unidad de medida correspondiente a cada cultivo (racimos para plátano, tachos o quintales para cacao), sin forzar una unidad única para todos los cultivos. \\
\hline
Métrica & Porcentaje de registros de cosecha con la unidad de medida correcta según el cultivo \\
\hline
Valor objetivo & 100~\%. \\
\hline
Método de medición & Verificación sobre una muestra de 30 registros de cosecha de al menos dos cultivos distintos. \\
\hline
Origen & EV-02 [02:08]; EV-01 [04:47]. \\
\hline
Prioridad & Alta \\
\hline
\end{longtable}
\end{center}

\bigskip
\noindent\textbf{RNF-08. Tiempo de generación de reportes de pérdidas}
\label{tab:RNF-08}

\begin{center}
\begin{longtable}{|p{4.2cm}|p{10cm}|}
\hline
Identificador & RNF-08 \\
\hline
Nombre & Tiempo de generación de reportes de pérdidas \\
\hline
Característica ISO/IEC 25010:2023 & Eficiencia de desempeño \\
\hline
Subcaracterística & Comportamiento temporal \\
\hline
Descripción & El sistema deberá generar el reporte de pérdidas por plagas (RF-19) dentro de un tiempo máximo, con un volumen representativo de registros cargados. \\
\hline
Métrica & Tiempo de respuesta \\
\hline
Valor objetivo & $\leq$ 3,0 s con 500 registros de enfermedad cargados. \\
\hline
Método de medición & Medición del percentil 95 de 30 ejecuciones. \\
\hline
Origen & ISO/IEC 25010:2023 (característica sin evidencia directa de campo). \\
\hline
Prioridad & Alta \\
\hline
\end{longtable}
\end{center}

\bigskip
\noindent\textbf{RNF-09. Tolerancia a interrupciones durante el registro}
\label{tab:RNF-09}

\begin{center}
\begin{longtable}{|p{4.2cm}|p{10cm}|}
\hline
Identificador & RNF-09 \\
\hline
Nombre & Tolerancia a interrupciones durante el registro \\
\hline
Característica ISO/IEC 25010:2023 & Fiabilidad \\
\hline
Subcaracterística & Tolerancia a fallos \\
\hline
Descripción & El sistema deberá conservar sin pérdida los datos ingresados en un formulario de registro (cosecha, labor o enfermedad) ante una interrupción de conexión durante el envío, permitiendo su reintento. \\
\hline
Métrica & Porcentaje de registros recuperados tras una interrupción simulada de conexión \\
\hline
Valor objetivo & 100~\%. \\
\hline
Método de medición & Prueba de interrupción de conexión simulada sobre 20 intentos de registro. \\
\hline
Origen & ISO/IEC 25010:2023, motivado por la pérdida de datos declarada en EV-01 [02:04], [02:17]. \\
\hline
Prioridad & Alta \\
\hline
\end{longtable}
\end{center}

\bigskip
\noindent\textbf{RNF-10. Bloqueo por intentos fallidos de acceso}
\label{tab:RNF-10}

\begin{center}
\begin{longtable}{|p{4.2cm}|p{10cm}|}
\hline
Identificador & RNF-10 \\
\hline
Nombre & Bloqueo por intentos fallidos de acceso \\
\hline
Característica ISO/IEC 25010:2023 & Seguridad \\
\hline
Subcaracterística & Resistencia \\
\hline
Descripción & El sistema deberá bloquear temporalmente una cuenta de usuario tras un número determinado de intentos fallidos consecutivos de inicio de sesión. \\
\hline
Métrica & Número de intentos fallidos permitidos antes del bloqueo, y duración del bloqueo \\
\hline
Valor objetivo & 5 intentos; bloqueo mínimo de 15 minutos. \\
\hline
Método de medición & Prueba de 5 intentos fallidos consecutivos sobre 10 cuentas de prueba. \\
\hline
Origen & CU-01 (Entrega 2, 1B); RD-04. \\
\hline
Prioridad & Alta \\
\hline
\end{longtable}
\end{center}

\bigskip
\noindent\textbf{RNF-11. Cifrado de credenciales almacenadas}
\label{tab:RNF-11}

\begin{center}
\begin{longtable}{|p{4.2cm}|p{10cm}|}
\hline
Identificador & RNF-11 \\
\hline
Nombre & Cifrado de credenciales almacenadas \\
\hline
Característica ISO/IEC 25010:2023 & Seguridad \\
\hline
Subcaracterística & Confidencialidad \\
\hline
Descripción & El sistema deberá almacenar las contraseñas de los usuarios mediante una función de cifrado unidireccional, sin conservar la contraseña en texto plano en ningún medio de almacenamiento. \\
\hline
Métrica & Porcentaje de contraseñas almacenadas en formato cifrado \\
\hline
Valor objetivo & 100~\%. \\
\hline
Método de medición & Revisión técnica del esquema de almacenamiento de credenciales en la base de datos. \\
\hline
Origen & ISO/IEC 25010:2023; R-04. \\
\hline
Prioridad & Alta \\
\hline
\end{longtable}
\end{center}

\bigskip
\noindent\textbf{RNF-12. Accesibilidad para usuarios sin experiencia digital previa}
\label{tab:RNF-12}

\begin{center}
\begin{longtable}{|p{4.2cm}|p{10cm}|}
\hline
Identificador & RNF-12 \\
\hline
Nombre & Accesibilidad para usuarios sin experiencia digital previa \\
\hline
Característica ISO/IEC 25010:2023 & Interacción de usuario \\
\hline
Subcaracterística & Inclusividad \\
\hline
Descripción & El sistema deberá permitir que un usuario sin experiencia previa en aplicaciones móviles complete el registro de una cosecha sin asistencia de un tercero. \\
\hline
Métrica & Porcentaje de usuarios de prueba sin experiencia previa que completan el registro sin asistencia \\
\hline
Valor objetivo & $\geq$ 80~\%. \\
\hline
Método de medición & Prueba de usabilidad con al menos 5 usuarios de prueba sin experiencia previa en aplicaciones. \\
\hline
Origen & EV-02 [00:24], [02:33], [02:43]. \\
\hline
Prioridad & Media \\
\hline
\end{longtable}
\end{center}

\bigskip
\noindent\textbf{RNF-13. Acceso desde dispositivo móvil}
\label{tab:RNF-13}

\begin{center}
\begin{longtable}{|p{4.2cm}|p{10cm}|}
\hline
Identificador & RNF-13 \\
\hline
Nombre & Acceso desde dispositivo móvil \\
\hline
Característica ISO/IEC 25010:2023 & Portabilidad \\
\hline
Subcaracterística & Adaptabilidad \\
\hline
Descripción & El sistema deberá ser accesible y operable desde un navegador móvil en un teléfono inteligente, sin requerir instalación de una aplicación nativa. \\
\hline
Métrica & Porcentaje de funciones principales operables desde un navegador móvil \\
\hline
Valor objetivo & 100~\% de las funciones de RF-01, RF-04 y RF-05. \\
\hline
Método de medición & Prueba funcional sobre un navegador móvil en al menos 2 tamaños de pantalla. \\
\hline
Origen & EV-01 [08:42]. \\
\hline
Prioridad & Alta \\
\hline
\end{longtable}
\end{center}

\bigskip
\noindent\textbf{RNF-14. Incorporación de nuevos cultivos sin modificar código}
\label{tab:RNF-14}

\begin{center}
\begin{longtable}{|p{4.2cm}|p{10cm}|}
\hline
Identificador & RNF-14 \\
\hline
Nombre & Incorporación de nuevos cultivos sin modificar código \\
\hline
Característica ISO/IEC 25010:2023 & Flexibilidad \\
\hline
Subcaracterística & Adaptabilidad \\
\hline
Descripción & El sistema deberá permitir agregar un nuevo tipo de cultivo y su unidad de medida asociada sin requerir modificaciones al código fuente de la aplicación. \\
\hline
Métrica & Tiempo requerido para incorporar un nuevo tipo de cultivo mediante configuración \\
\hline
Valor objetivo & $\leq$ 30 minutos, sin intervención de un desarrollador. \\
\hline
Método de medición & Prueba de incorporación de un cultivo de prueba por parte de un administrador capacitado. \\
\hline
Origen & ISO/IEC 25010:2023; RD-05. \\
\hline
Prioridad & Media \\
\hline
\end{longtable}
\end{center}

\bigskip
\noindent RNF-15, correspondiente a la explicabilidad de las alertas y recomendaciones del módulo de IA, se documenta íntegramente en la Sección~\ref{sec:explicabilidad-ia}, ficha de la Tabla~\ref{tab:RNF-15}.

% Bibliografia
\bibliographystyle{ieeetr}
\bibliography{referencias}

\end{document}
